% Generated draft: requires mathematical and source review.
\documentclass[11pt]{article}
\usepackage[margin=25mm]{geometry}
\usepackage{fontspec}
\setmainfont{Latin Modern Roman}
\usepackage{amsmath,amssymb,mathtools,unicode-math}
\setmathfont{Latin Modern Math}
\usepackage{enumitem,needspace,xurl,hyperref}
\setlength{\parindent}{0pt}
\setlength{\parskip}{5pt}
\setlength{\emergencystretch}{3em}
\newcommand{\N}{\mathbb{N}}
\newcommand{\Z}{\mathbb{Z}}
\newcommand{\Q}{\mathbb{Q}}
\newcommand{\R}{\mathbb{R}}
\newcommand{\C}{\mathbb{C}}
\newcommand{\F}{\mathbb{F}}
\newcommand{\A}{\mathbb{A}}
\newcommand{\PP}{\mathbb{P}}
\newcommand{\E}{\mathbb{E}}
\newcommand{\eps}{\varepsilon}
\newcommand{\dd}{\,\mathrm{d}}
\begin{document}
\section*{UnsolvedMath ID: 30005494}
\textbf{Source problem code:} catalogue entry
\subsection{Definitions and mathematical statement}
A complex algebraic singularity germ \((X,x)\) is called rational if \(X\) is normal near \(x\) and, for a resolution of singularities \(f\colon Y\to X\), one has \(R^i f_*\mathcal O_Y=0\) near \(x\) for every \(i>0\).

For the topological link convention, choose a local embedding \((X,x)\hookrightarrow(\C^N,0)\) and a sufficiently small Euclidean sphere \(S^{2N-1}_\varepsilon\). Set
\[
L(X,x)=X\cap S^{2N-1}_\varepsilon,
\qquad
\pi_1^{\mathrm{link}}(X,x)=\pi_1(L(X,x),\ell),
\]
where \(\ell\in L(X,x)\) is a base point. The group is considered up to isomorphism. The supplied record does not explicitly define its notation \(\pi_1(X;x)\). In particular, it does not settle whether the intended group is this link group or the fundamental group of the smooth locus of a sufficiently small neighborhood. These constructions must not be identified without appropriate hypotheses.

\textbf{Classification question, conditional on the link convention.} Characterize the class
\[
\mathcal R_{\mathrm{link}}=
\left\{[G]\;\middle|\;
\begin{gathered}
G\text{ is a finitely presented group, and there exists a rational}\\
\text{complex algebraic singularity germ }(X,x)\text{ such that}\\
G\cong\pi_1^{\mathrm{link}}(X,x)
\end{gathered}
\right\}.
\]
Equivalently, determine necessary and sufficient conditions on a finitely presented group \(G\) for such a germ to exist. The clean statement imposes neither a fixed complex dimension nor an isolated-singularity hypothesis; neither restriction is added here. The precise local-fundamental-group convention remains a source issue requiring review.

\textbf{Supplied surrounding results and qualifications.} For a group \(G\), write \(H_i(G;\Q)\) for group homology with trivial rational coefficients. The group is \(\Q\)-perfect if \(H_1(G;\Q)=0\), and \(\Q\)-superperfect if \(H_1(G;\Q)=H_2(G;\Q)=0\). For a resolution with a simple normal crossings exceptional divisor, its dual complex records the exceptional components and their intersections; denote it by \(D(X;x)\).

The supplied excerpt labels the following as Theorem 7: for a rational singularity \((X,x)\), the group \(\pi_1(D(X;x))\) is \(\Q\)-superperfect; moreover, for every finitely presented \(\Q\)-superperfect group \(G\), there exists a rational singularity \((X_G,x)\) of complex dimension \(6\) with \(\pi_1(X_G;x)\cong G\). These are recorded as supplied claims. The necessity assertion concerns the dual-complex group, whereas the existence assertion uses the local-group notation. The excerpt does not justify transferring the necessity assertion to the local group, and its resolution hypotheses need checking in the full source.

The supplied Theorem 8 states that, for every finitely presented group \(G\), being \(\Q\)-perfect is equivalent to being the fundamental group of a Cohen--Macaulay singularity, attributed to Koll\'ar, reference [10, Theorem 4]. The excerpt does not specify all conventions or additional hypotheses for this separate class of singularities.

The dated curation also reports complete realizability for normal isolated threefold singularities, while describing the rational and Cohen--Macaulay classifications collectively as partial. It does not state that the threefold realizability result imposes rationality. Its description of the Cohen--Macaulay case requires reconciliation with the excerpt's claim of a complete answer in that case.

For context, supplied Theorem 9 asserts that there is a constant \(c_1(n)\), depending only on \(n\), such that every \(n\)-dimensional log terminal singularity has a local fundamental group admitting a finite normal abelian subgroup \(A\) with
\[
\operatorname{rank}(A)\le n,
\qquad
[\pi_1(X;x):A]\le c_1(n).
\]
Here the rank of a finite abelian group means its minimum number of generators. The excerpt attributes this result to Braun, Filipazzi, Svaldi, and the unnamed report author. It says that the proof uses local-to-global methods from the Minimal Model Program and Birkar's boundedness of singular Fano varieties. It further says that \(c_1(n)\) is generally not the constant \(c_0(n)\) in its Theorem 5 and that sharp bounds for \(c_1(n)\) were unclear. A subsequent assertion bounding \(\operatorname{rank}(A)\) by a regularity invariant is truncated, so no complete statement of that assertion is available. These contextual claims do not supply a classification for the question above.
\subsection{Short English statement}
Determine exactly which finitely presented groups arise as local fundamental groups of rational complex algebraic singularities, with dimension unrestricted and no isolatedness assumption added. The supplied excerpt distinguishes local groups from groups of dual complexes, but leaves the local-group convention and some surrounding hypotheses unresolved. This is a statement-only draft requiring human mathematical review.
\subsection{Sources}
\begin{itemize}
\item Local registry: \texttt{data/unsolvedmath/problems.json}, with \texttt{catalog/problems.json} as the supplementary registry; supplied record for permanent UnsolvedMath ID \(30005494\), title \emph{Local Fundamental Groups of Rational Singularities}, exact source code \texttt{catalogue entry}. These registry paths identify provenance; their contents have not been independently inspected here.
\item Supplied authoritative report citation: \emph{Komplexe Analysis -- Differential and Algebraic methods in Kähler spaces} (2023), Oberwolfach Reports, pp.\ 965--1030, DOI \texttt{10.4171/owr/2023/18}. Supplied source link: \url{https://doi.org/10.4171/owr/2023/18}. The extracted text ends on report page 1007. No proposer is named in the record.
\item Supplied literature source, labelled ``Kapovich and Kollar'': \url{https://doi.org/10.1090/S0894-0347-2014-00807-9}. The record supplies no article title or full bibliographic details. The surrounding text attributes work on dual-complex fundamental groups to Koll\'ar and Kapovich.
\item Supplied literature source, labelled ``OWR 2023 report'': \url{https://ems.press/content/serial-article-files/47013}.
\item Literature assessment checked at \texttt{2026-08-22}: the supplied status is \texttt{partially\_solved}. The supplied assessment reports complete realizability for normal isolated threefold singularities and partial rational and Cohen--Macaulay classifications, without a complete answer for all singularities in the question. This is supplied dated curation, not an independently verified literature assessment; the record expressly requests independent verification of time-sensitive status and cited claims.
\item The registry labels the statement \texttt{corrected\_verified} and reports that it was checked against the official Oberwolfach report. This is supplied verification metadata, not a verification performed for this draft.
\item Internal references in the excerpt are incomplete: [10, Theorem 4] is attributed to Koll\'ar; [2, Theorem 2] to Braun, Filipazzi, Svaldi, and the unnamed report author; [1, Theorem 1.1] to Birkar; and [13, Theorem 7] to the unnamed report author. The excerpt also refers to its Theorem 5 without supplying that statement. No bibliographic entries or additional URLs for these numbered references are supplied.
\end{itemize}
\end{document}
